\begin{abstract}
The tidal screen of the companion paper, with the Lagrangian constant set to zero, puts new comets on the wrong side of the observed Oort peak. A comet that falls in from rest arrives with a Newtonian \(1/a\) of about \(1.23\times10^{-4}\,\mathrm{au^{-1}}\), so its original semimajor axis is at most about \(8{,}100\,\mathrm{au}\). The observed peak of new comets lies between \(10^{4}\) and \(10^{5}\,\mathrm{au}\). The push sits in a shell near \(9\) to \(12\) thousand au, where the net gravity on a test body points outward. Three evaluations of the same map agree at the percent level. A radial orbit integration at the same \(C=0\) lands the same way: starts of \(3\times10^{4}\) and \(9.5\times10^{4}\,\mathrm{au}\) arrive at \(4{,}877\) and \(6{,}270\,\mathrm{au}\), and two starts inside the shell never reach the inner system. This paper records that null. It does not introduce a new constant.
\end{abstract}

\section{What was frozen}

The calculation was specified before the screened numbers were computed. The map uses radii from \(0.05\) to \(300\,\mathrm{kAU}\) and twelve Gauss--Legendre nodes in the cosine of the angle from the external field. The external-field strengths are \(1.9\times10^{-10}\) and \(2.32\times10^{-10}\,\mathrm{m\,s^{-2}}\). The action-form screen is the one whose \(x^*\) was fixed by Cassini at \(7.04\times10^{11}\). The constant \(C\) in the shell stress is zero, so the stress is proportional to \(P(y)\), not to \(P(y)-C\).

A comet falling from rest at infinity has an osculating Newtonian binding
\[
\frac{1}{a}=\frac{2\Phi(0)}{GM_\odot},
\]
with \(\Phi(0)\) the line integral of the anomalous acceleration. Energy conservation fixes this number. The torque along the orbit is a separate statement and is not the kill.

\section{The binding offset}

At \(g_e=1.9\times10^{-10}\,\mathrm{m\,s^{-2}}\) the offset is \(+1.23\times10^{-4}\,\mathrm{au^{-1}}\). At \(2.32\times10^{-10}\) it is \(+1.61\times10^{-4}\,\mathrm{au^{-1}}\). Every comet that the model places on a bound orbit from rest therefore has an original semimajor axis of at most about \(8{,}100\,\mathrm{au}\). The observed peak lies at \(10^{4}\)--\(10^{5}\,\mathrm{au}\), so the model leaves that peak empty.

\begin{center}
\begin{tabular}{lcc}
\hline
evaluation & \(1.9\times10^{-10}\) & \(2.32\times10^{-10}\) \\
\hline
line integral, 12 directions & \(+1.23\times10^{-4}\) & \(+1.61\times10^{-4}\) \\
phantom-source potential at the Sun & \(+1.227\times10^{-4}\) & \(+1.610\times10^{-4}\) \\
\hline
\end{tabular}
\end{center}

The twelve directions agree with each other at the level of \(0.1\) percent. The phantom-source split is a divergence piece of \(-1.9\times10^{-5}\) and a shell piece of \(+1.42\times10^{-4}\). The shell raises the potential inside itself. A separate \(W\)-only radial integral, without the divergence piece, gives \(1.386\times10^{-4}\,\mathrm{au^{-1}}\) at \(1.9\times10^{-10}\), and an edge at \(7{,}216\,\mathrm{au}\). The two calculations are the same shell. They are not the same integral. The number that kills \(C=0\) here is the full map, \(1.23\times10^{-4}\,\mathrm{au^{-1}}\).

\section{Where the gravity points outward}

The orientation-averaged anomalous acceleration, divided by the solar Newtonian acceleration, reaches \(+0.97\) at \(8.7\,\mathrm{kAU}\), \(+1.27\) at \(10\,\mathrm{kAU}\), and \(+0.99\) at \(11.5\,\mathrm{kAU}\). Between about \(8.7\) and \(11.5\,\mathrm{kAU}\) the net radial gravity on a test body points outward. A body placed there is not bound to the Sun in the usual sense.

Radial integrations at \(C=0\), with angular momentum aimed at a \(5\,\mathrm{au}\) perihelion, give the same conclusion orbit by orbit. A start at \(1.39\times10^{4}\,\mathrm{au}\) turns at \(9{,}683\,\mathrm{au}\) and does not reach the inner system. A start at \(3.0\times10^{4}\,\mathrm{au}\) reaches perihelion with osculating semimajor axis \(4{,}877\,\mathrm{au}\). A start at \(9.49\times10^{4}\,\mathrm{au}\) arrives at \(6{,}270\,\mathrm{au}\). Both arrivals sit below \(10^{4}\,\mathrm{au}\). The osculating \(1/a\) matches the potential to \(1.3\times10^{-6}\,\mathrm{au^{-1}}\) on the orbits that do reach perihelion.

\section{Why a larger screen radius is not available}

The offset scales as \(1/x^*\). Cassini already sets the largest \(x^*\) the action form is allowed, \(7.04\times10^{11}\). Moving \(x^*\) up to weaken the Oort push moves the Cassini quadrupole through its bound. There is no remaining \(x^*\) that keeps both. This paper does not retune \(a_0\) or the interpolating function.

The torque proxy is consistent with the same shell and is not the kill. The jump semimajor axis is \(3.67\,\mathrm{kAU}\) and \(4.33\,\mathrm{kAU}\) at the two external fields, against \(2.63\,\mathrm{kAU}\) with the screen removed and \(26.9\,\mathrm{kAU}\) for Newton. Newton's jump falling between \(15\) and \(40\,\mathrm{kAU}\) was the check that the integrator is the one that was pre-registered.

\section{What remains open}

A population synthesis of the comet catalogue, with planetary perturbations and the Galactic tide acting together over the age of the Solar System, has not been run at \(C=0\). The statement here is the energy integral, checked by direct radial orbits. Calling the observed peak empty is the comparison of that integral with the published range \(10^{4}\)--\(10^{5}\,\mathrm{au}\). The population synthesis is open.

The companion paper studies the same equation at the nonzero \(C\) that moves the edge into the observed window. That nonzero value is a different model. This paper does not adopt it. It records that \(C=0\) fails.

